% HISTORICAL DRAFT -- NOT PART OF THE MANUSCRIPT OR CURRENT EVIDENCE.
% This file is not included by paper/main.tex and contains unverified citations
% and superseded LOCUS mechanism descriptions. Do not input, cite, or ship it as
% submission material unless every claim and reference is reconciled against
% verified primary sources. The authoritative comparison is maintained in
% docs/related-work-comparison.md and the Related Work section of main.tex.
\section{Related Work}
\label{sec:related}

Private-key backup and recovery spans threshold secret sharing, privacy-preserving password-based retrieval, social recovery, biometric or memory-based reconstruction, and hardware-assisted recovery. We focus on the work most relevant to \textsf{LOCUS}: distributed trust, human-centered recovery inputs, privacy-preserving authorization, and verifiable share return.

\subsection{Threshold Recovery and Verifiable Secret Sharing}

Threshold cryptography provides the basic mechanism for avoiding a single point of failure in key protection and recovery. Shamir's secret sharing established the standard \((t,n)\) model, while later work on verifiable secret sharing, publicly verifiable secret sharing, and proactive resharing extended it to settings with faulty participants, public verification, and share refresh or recovery \cite{Shamir1979,Feldman1987,PVSS1999,AVSS2002}. More recent systems have applied threshold recovery to blockchains, cloud key management, threshold signatures, rational secret sharing, timed release, and dynamic storage environments \cite{ConsortiumBlockchain2021,CloudKMS2021,ThresholdECDSARecovery2022,RandomBeaconRecovery2022,HierarchicalStorage2023,RationalSecretSharing2024,PRTDR2024}.

\textsf{LOCUS} follows this distributed-trust direction but differs in two ways. First, it does not secret-share the user's private key \(sk_U\) directly. Instead, it samples a random recovery key \(K_R\), wraps \(sk_U\) under a symmetric key derived from \(K_R\), and shares only \(K_R\). Second, share release is tied to a fresh authorization flow derived from the user's remembered recovery input. In this sense, \textsf{LOCUS} combines threshold protection with an explicit online authorization layer, whereas much of the threshold-recovery literature focuses primarily on share distribution, refresh, or reconstruction.

The protocol is also close in spirit to blockchain-oriented recovery schemes that use verifiable secret sharing \cite{ConsortiumBlockchain2021,SocialWallet2023}. Those schemes typically rely on certificates, smart contracts, or guardian identities as the release condition. By contrast, \textsf{LOCUS} uses threshold sharing only for custody of \(K_R\), while authorization is driven by a privacy-preserving function of a user-reproducible recovery input.

\subsection{Password-Based and Privacy-Preserving Retrieval}

A closely related line of work studies recovery from a human-memorizable secret while reducing server-side exposure through oblivious cryptographic techniques. Password-protected and password-authenticated secret-sharing schemes show how multiple servers can help a user recover a secret using only a password while preventing below-threshold or otherwise unauthorized views from becoming offline password-testing oracles \cite{Bagherzandi2011PPSS,Jarecki2014PPSS,PPSS2016,Camenisch2014Memento,Camenisch2015TPASS,Yi2015TPASS,Yi2019TPASS}. More recent Password-Protected Key Retrieval (PPKR) designs refine this model and explore different trust assumptions, including settings with and without HSM protection \cite{PPKR2024}. Privacy-preserving account-recovery systems have also used OPRF-like primitives to validate recovery information without fully exposing it to the service provider \cite{SecureAccountRecovery2024}. The security analysis of the WhatsApp encrypted backup protocol similarly highlights the practical importance of challenge-response authorization and server-side rate limiting when recovery relies on a human-chosen secret \cite{WhatsAppBackup2023}.

\textsf{LOCUS} is closest to this family in its use of a privacy-preserving authorization layer. The difference is that our protocol uses a deterministic encoding of a deployment-defined recovery input instead of a conventional password or security-question answer. That input is hashed into the cue-derived password used by TPASS. Moreover, the human-reproducible input is not itself the recovered private key; it authorizes threshold recovery of wrapping material for a separately encrypted key. This separation between \emph{authorization} and \emph{reconstruction} distinguishes \textsf{LOCUS} from password-protected retrieval systems in which the human-memorizable secret is more directly tied to key reconstruction.

\subsection{Authentication, Password Managers, and Account Recovery}

Web authentication and recovery remain unsettled as systems move toward passkeys, end-to-end encrypted accounts, password managers, and platform sync services \cite{Blessing2025WebAuthRecovery,Barbosa2025FIDO2,Li2014PasswordManagers,Pearman2019PasswordManagers,FIDOAlliancePasskeys,AppleKeychainSecurity}. These mechanisms improve day-to-day credential availability, but recovery authority is typically mediated by a provider account, device ecosystem, browser, password-manager vault, or account-recovery workflow.

\textsf{LOCUS} addresses a complementary problem: recovering a user-controlled private key after device loss without introducing a password verifier, centralized escrow service, or social helper who can reconstruct the key. It is therefore closer to recovery for encrypted credentials than to daily authentication.

\subsection{Social Recovery and Community-Assisted Schemes}

Social recovery schemes replace or supplement cryptographic custody with interpersonal trust. Early proposals used prior transaction partners or arbitration processes to recover decentralized accounts \cite{DAppRecovery2019}. Later systems employed hidden assistance relationships with zero-knowledge proofs \cite{CRSA2021}, guardian-based smart contracts \cite{SocialWallet2023}, community-based key reconstruction models such as Bottom-Up Secret Sharing \cite{ANARKey2025}, and peer-assisted passkey recovery in enterprise settings \cite{PeerRecovery2025}.

These works address the same high-level problem as \textsf{LOCUS}: recovery without a single central authority. Their main recovery factor, however, is \emph{social authorization}. Recovery succeeds because a threshold of guardians, assistants, peers, or community members attests to the user or contributes recovery material. \textsf{LOCUS} instead uses a user-held, privacy-preserved recovery input derived from a flexible cue policy. The parties are still distributed and threshold-based, but they do not act as social witnesses of identity. This avoids some social dependencies, such as guardian selection and relationship stability, while shifting more weight onto the entropy and reproducibility of the user's chosen cues.

\subsection{Biometric, Memory-Based, and Human-Centered Recovery Factors}

Prior work shows that human-chosen passwords and personal-knowledge answers can be memorable yet biased, predictable, and exposed through public or social information \cite{Bonneau2012Quest,Bonneau2015Secrets,Yan2004Password}. Graphical and cue-based authentication similarly explores the tradeoff between recall and guessability \cite{Biddle2012Graphical}. Another major line of work derives or reconstructs cryptographic material from noisy human or physical inputs. Fuzzy commitments and fuzzy extractors recover stable secrets from biometric measurements \cite{FuzzyCommitment1999,FuzzyExtractors2004}. Later works explored PUF-based recovery \cite{PUF2011,RBC2020}, iris-based authentication \cite{Iris2019}, DNA-based recovery \cite{DNABioCrypto2023}, biometric protection of mnemonic phrases \cite{LeakableMnemonic2024}, and multimodal systems that combine autobiographical memory with biometrics \cite{LifeXP2025}.

These approaches share with \textsf{LOCUS} an interest in moving beyond stored digital backup secrets toward more human-centered recovery factors, but the technical model is different. Biometric and PUF-based systems must tolerate noisy inputs and therefore rely on fuzzy matching, helper data, or pattern reconstruction. By contrast, \textsf{LOCUS} uses deterministic cue policies that map selected recovery inputs to stable identifiers under public canonicalization rules. Location--Person cues are one prototype policy, motivated by prior work on contextual and cue-based recall, but they are not the paper's contribution. In this sense, the protocol is closer to structured memory-based recovery than to biometric extraction. The nearest comparison is Reminisce \cite{Reminisce2022}, which also uses personal memory and location-linked information for blockchain key recovery. However, Reminisce uses picture-based personal memory to regenerate a temporary key, whereas \textsf{LOCUS} uses cue-derived input only to authorize threshold recovery of wrapping material.

\subsection{Hardware-Assisted and Custodial Recovery Frameworks}

A final line of work uses secure hardware, TEEs, custodial infrastructures, or auditable recovery services to protect secrets and mediate recovery \cite{Omega1996,Acsesor2022,WhatsAppBackup2023,CloudRecovery2025,MFAKeystore2025}. Related key-lifecycle work studies post-compromise secure cloud storage, continuous key rotation, and secure-element mechanisms for slowing password guesses or supporting deniable encryption \cite{Gloudemans2025Akeso,Lim2024QUICKeR,Hugenroth2024Sloth}. These systems often provide stronger server-side control over brute-force attempts, transparent logging, or flexible policy enforcement, and some can support recovery without exposing the protected secret outside a secure module.

\textsf{LOCUS} differs from these approaches in that it does not assume a TEE, HSM, or centralized recovery authority as a core component. Instead, it pushes release logic into a threshold set of independent parties and makes one-time challenges, expiration, lockout, and logging explicit protocol-level requirements. This gives the design a more decentralized profile, although some protections offered naturally by hardware-backed systems remain operational assumptions here.

\subsection{Usability and Resolver Privacy}

Recent privacy-enhancing-technology studies emphasize that subtle cryptographic features need empirical validation because users' perceptions, disclosure decisions, and support needs can shape adoption \cite{Fischer2025KeyTransparency,Dehling2024BiometricDisclosure,Demuth2024SupportPersonas}. Contextual cryptographic workflows have also been explored with environmental inputs from IoT devices \cite{Jois2024SocIoTy}.

\textsf{LOCUS} similarly uses human-reproducible context, but with a distinct privacy boundary. Cue inputs are processed locally into a threshold password input, while the baseline system does not hide lookup activity from resolver providers. This separates storage-recovery privacy from resolver privacy and motivates the paper's explicit treatment of resolver-observation attackers.

Overall, the closest prior work to \textsf{LOCUS} comes from three directions: threshold-based key recovery, privacy-preserving password-style retrieval, and human-centered recovery mechanisms. Our protocol combines elements from each, but differs in one central design decision: it uses a deterministic cue-derived factor only to authorize recovery privately, while placing actual key restoration in a separate threshold recovery layer.
